\documentclass[11pt]{article}
\usepackage{amsmath, amssymb}
\usepackage[margin=0.85in]{geometry}
\usepackage{booktabs}
\usepackage{siunitx}
\usepackage{longtable}
\usepackage{array}
\usepackage{seqsplit}
\usepackage{hyperref}

\DeclareSIUnit{\usd}{USD}
\DeclareSIUnit{\year}{yr}
\newcolumntype{P}[1]{>{\raggedright\arraybackslash}p{#1}}
\newcommand{\code}[1]{\texttt{\seqsplit{#1}}}

\title{Model Documentation: Governing Equations and Economics for a\\
Solar-Driven, Single-Bed PAM--LiCl SAWH Device}
\author{}
\date{}

\begin{document}
\maketitle

\section*{Reference and scope}

This document describes the lumped device + techno-economic model implemented in
\code{solar\_lumped}, following Wilson \& D\'{i}az-Mar\'{i}n \textit{Device} (2025) Eqs.~1--6
and Note~S1 (COMSOL Table~S3 parameters) for the physics, and the package's own
\code{economics/} module for cost. The device is a single hydrogel bed that alternates
between open \emph{absorption} (night) and sealed \emph{desorption} (day) half-cycles. This
is the default (\code{desorption\_solver="quasi\_steady"}) configuration; one alternate
exists in code but is out of scope here: \code{segregated} / \code{coupled\_bdf} desorption
integrators, which trade the quasi-steady algebraic thermal solve for COMSOL-mimicking or
fully transient alternatives (\code{device\_config.py}).

\textbf{Base case is Case 2, not Wilson's original paper physics.} As of 2026-07,
\code{DeviceConfig.thermal\_params()} -- the physics every plain \code{DeviceConfig()} call
gets unless it explicitly overrides \code{thermal=} -- defaults to real absorber/glass IR
emissivities ($\varepsilon_{abs,ir}=0.05$, $\varepsilon_{gl,ir}=0.95$, the "selective surface"
variant, Case~2 in \code{sawh\_bayesopt.design\_space.CASE\_EPS\_IR}) rather than Wilson's
original blackbody/cavity approximation (Case~1). Every equation, assumption, and default value
below describes Case~2 unless stated otherwise. Case~1 (Wilson's exact paper physics) remains
reachable via an explicit \code{DeviceThermalParams(eps\_abs\_ir=1.0, eps\_glass\_ir=1.0)} (or
the equivalent \code{eps\_abs\_ir=eps\_glass\_ir=None}, still numerically identical -- see
Section~\ref{sec:radiative}) override, and is exactly what the paper-reproduction scripts under
\code{analysis/paper\_recreation/wilson/} use -- those scripts build their own
\code{DeviceThermalParams} directly rather than going through
\code{DeviceConfig.thermal\_params()}, so they are unaffected by this default and continue to
reproduce Wilson's digitized figures under Case~1. Case~3 (idealized optical-material limits,
$\varepsilon_{abs,ir}=\varepsilon_{gl,ir}=0.0$) is a third explicit sensitivity variant, also not
the default.

\textbf{Subscripts:} $g$ = hydrogel; $abs$ = solar absorber; $gl$ = cover glass; $cond$ =
aluminum condenser; $amb$ = ambient.

\textbf{How to use the Source columns below:} cells already filled cite the in-repo file or
paper section the value/range comes from. Cells marked \textbf{TBD} have no literature or
vendor citation on file yet -- the value is either an as-implemented constant or an
arbitrarily chosen sweep bound with no documented justification. Fill these in as sources are
found; do not treat a TBD as validated.

\tableofcontents

\section{Model design}

\subsection{Device overview}

A footprint-area device stacks (top to bottom, day-side up): cover glass $\to$ air/insulation
gap $\to$ solar absorber (aluminum) $\to$ hydrogel-salt composite bed $\to$ vapor gap $\to$
aluminum condenser (finned, ambient-cooled). All balances are per unit footprint area
(\si{\metre\squared}) unless noted. During \textbf{desorption} (daytime), the chamber is
sealed: solar radiation heats the absorber, which conducts heat into the gel; the gel desorbs
water vapor across the vapor gap to the cooler condenser, where it condenses and is collected.
During \textbf{absorption} (nighttime), the chamber is opened to ambient air; the gel re-hydrates
by uptake from ambient humidity, with no solar input and no condenser dynamics. An optional
uncovered variant (\code{has\_glass=False}) removes the glass/insulation-gap stack entirely and
radiates the absorber directly to a sky/ambient sink.

\subsection{State variables}

\begin{table}[h]
\centering
\small
\begin{tabular}{@{}lll@{}}
\toprule
\textbf{Symbol} & \textbf{Units} & \textbf{Role} \\
\midrule
$c_w$ & \si{\mol\per\cubic\metre} & Brine water concentration in gel (Eq.~5 state) \\
$H$ & \si{\metre} & Hydrogel thickness (Eq.~6 state; reference $H_0$) \\
$T_{cond}$ & \si{\celsius} & Condenser temperature (transient during desorption only) \\
$T_g, T_{abs}, T_{gl}$ & \si{\celsius} & Quasi-steady thermal states (algebraic each step) \\
\bottomrule
\end{tabular}
\end{table}

\textbf{Absorption phase ODEs:} $\dot{c}_w$, $\dot{H}$ with $Q_{solar}=0$, $\dot{m}_{des}=0$,
$\dot{T}_{cond}=0$. \textbf{Desorption phase ODEs:} $\dot{c}_w$, $\dot{H}$, $\dot{T}_{cond}$
with $T_g$, $T_{abs}$, $T_{gl}$ from the coupled thermal solve.

\subsection{Daily cycle structure}

One simulated day is one absorption half-cycle followed by one desorption half-cycle, gel
state carried forward between them. The absorption/desorption split follows the forcing
weather profile's own real day/night boundary (solar-radiation threshold crossing, so the two
halves are not forced to equal length) in \code{real}/replay weather modes; the
\code{baseline} synthetic-profile mode instead defaults to a fixed
\SI{12}{\hour}/\SI{12}{\hour} split (\code{weather/profiles.py::PHASE\_HOURS}) so that
single-axis parameter sweeps hold the forcing schedule constant. Multi-day runs either warm up
for a fixed cycle count or use Aitken $\Delta^2$ extrapolation (\code{find\_cyclic\_state}) to
locate the true steady periodic post-desorption $(c_w, H)$ state before reporting a single
representative day.

\section{Physics model assumptions}

\begin{itemize}
\item \textbf{Lumped / 0-D:} the gel, absorber, glass, and condenser are each a single uniform
temperature node per footprint area; no through-thickness gradients are resolved.
\item \textbf{Quasi-steady thermal solve during desorption:} $T_g$, $T_{abs}$, $T_{gl}$ are found
by root-finding the steady residual form of Eqs.~1, 3, 4 at every ODE evaluation (assumes the
thermal response is fast relative to the hydration/mass-transfer timescale). Only $T_{cond}$
carries real thermal inertia (finite aluminum slab thermal mass) during desorption.
\item \textbf{Absorption thermal shortcut:} during absorption, $T_g = T_{abs} = T_{gl} = T_{amb}$
is assumed outright (fast gel thermal storage, \code{coupled\_dynamics.py} Note~S1 Eq.~S1)
rather than solved from Eqs.~1/3/4 -- there is no solar input and no condenser coupling to
balance against.
\item \textbf{Radiative exchange (base case, Case~2 "selective surface"):}
absorber$\leftrightarrow$glass and glass$\leftrightarrow$ambient exchange use real IR
emissivities $\varepsilon_{abs,ir}=0.05$, $\varepsilon_{gl,ir}=0.95$ (Table~\ref{tab:phys-swept}),
not Wilson's typeset Eqs.~3--4 blackbody/cavity approximation ($\varepsilon=1$). Setting both to
$1.0$ (Case~1) reproduces Wilson's original paper physics exactly; setting both to $0.0$ (Case~3)
gives the idealized optical-material-limits sensitivity variant. See
Section~\ref{sec:radiative}.
\item \textbf{Radiative sink temperature:} the default sink for the outermost surface is ambient
air temperature; an optional \code{sky\_temp\_depression\_c} lowers the effective radiant sky
temperature below air temperature, needed to reproduce the COMSOL Atacama field-test curves
(surface-to-ambient radiation with a real sky, not blackbody air).
\item \textbf{Water inventory basis:} $c_w$ and $H$ are tracked as separate states, but the
sorbate mass used for activity/uptake calculations is referenced to the fixed fabrication
thickness $H_0$, not the swollen $H(t)$ -- swelling only feeds back into vapor-gap convection
($L_g-H$) and gel--absorber conductance ($H/k_{gel}$). This avoids double-counting dilution (see
\code{salt\_properties.py::pam\_licl\_gravimetric\_uptake\_g\_g}).
\item \textbf{Water activity (LiCl):} brine activity from the Conde (2004) correlation, capped
during absorption only by the measured PAM--LiCl DVS isotherm, i.e.\ $a_w =
\max(a_{w,brine}, a_{w,DVS})$ in absorption; brine activity alone in desorption. Non-LiCl salts
(NaCl, CaCl$_2$, MgCl$_2$) use brine activity only, with no DVS cap (no isotherm digitized for
them) -- out of this document's primary scope but present in code.
\item \textbf{Mass-transfer coefficient:} absorption uses a fixed empirical open-chamber
convection coefficient $g_{chamber}$; desorption derives $g$ from the Hollands natural-convection
$h_{conv,g}$ via a Lewis-number-unity heat/mass analogy (Note~S1 Eq.~S5). No forced convection is
modeled in either phase.
\item \textbf{Self-consistent desorption coupling:} $\dot{m}_{des}$ is found as the root of
$\dot{m}_{calc}(\dot{m}_{des}) - \dot{m}_{des} = 0$ (Brent's method) so that the gel heat balance
(Eq.~1) and the mass-transfer equations (Eqs.~5--6) agree on the same rate at every step.
\item \textbf{Scope:} default salt is LiCl with a hydrogel sorbent. A MOF-sorbent path and three
additional deliquescent salts exist in code with their own property tables but are not documented
further here.
\item \textbf{Initial condition:} $c_w$ initializes at equilibrium with the $\sim$\SI{20}{\percent}
RH ambient the gel is fabricated at (DVS isotherm, LiCl), unless a cyclic steady-state or explicit
override is requested instead.
\end{itemize}

\section{Thermal balances (Wilson Eqs.~1, 3, 4)}\label{sec:radiative}

During each integration step the gel, absorber, and glass temperatures satisfy the steady
residual form (Eqs.~1, 3, 4 set to zero):

\begin{equation}
U_{gel}(T_{abs} - T_g)
= h_{conv,g}(T_g - T_{cond})
+ \dot{m}_{des}\,h_{des}
+ \varepsilon_{gc}\sigma(T_g^4 - T_{cond}^4)
\label{eq:gel}
\end{equation}

\begin{equation}
k_{air}\frac{T_{abs} - T_{gl}}{L_{ins}}
+ \varepsilon_{ag}\sigma(T_{abs}^4 - T_{gl}^4)
= h_{amb}(T_{gl} - T_{amb})
+ \varepsilon_{ga}\sigma(T_{gl}^4 - T_{sky}^4)
\label{eq:glass}
\end{equation}

\begin{equation}
\varepsilon_{abs}\,\tau_{glass}\,Q_{solar}
= k_{air}\frac{T_{abs} - T_{gl}}{L_{ins}}
+ \varepsilon_{ag}\sigma(T_{abs}^4 - T_{gl}^4)
+ U_{gel}(T_{abs} - T_g)
\label{eq:absorber}
\end{equation}

where $T_{sky} = T_{amb} - \Delta T_{sky}$ (default $\Delta T_{sky}=0$). In the uncovered variant
(\code{has\_glass=False}), Eqs.~\eqref{eq:glass}--\eqref{eq:absorber} collapse to a single
absorber-to-sky balance $\varepsilon_{abs}Q_{solar} = h_{amb}(T_{abs}-T_{amb}) +
\varepsilon_{abs}\sigma(T_{abs}^4 - T_{sky}^4) + U_{gel}(T_{abs}-T_g)$ with $T_{gl} \equiv T_{amb}$.

\textbf{Correlations (Note~S1):}
\begin{itemize}
\item $h_{conv,g}$: Hollands et al.\ natural convection between parallel plates in the vapor gap
$L_g - H$ (tilt-dependent; Eqs.~S3--S4).
\item $\varepsilon_{gc} = 1/(1/\varepsilon_{gel} + 1/\varepsilon_{Al} - 1)$: parallel-plate
radiative exchange (Eq.~S2).
\item $U_{gel}$: series-resistance gel--absorber conductance (updated from the single-layer
approximation in earlier drafts of this document to match \code{table\_s3.py::u\_gel\_w\_m2\_k}):
\begin{equation}
\frac{1}{U_{gel}} = \frac{L_{Al}}{k_{Al}} + \frac{L_{silicone}}{k_{silicone}} + \frac{H(t)}{k_{gel}}
\label{eq:ugel}
\end{equation}
(aluminum backing + silicone bonding layer + hydrogel conduction in series; $H$ floored at
$H_0/4$ to avoid a singular resistance).
\item $\varepsilon_{ag}, \varepsilon_{ga}$ (base case, Case~2): $\varepsilon_{ag} =
\mathrm{parallel\_plate\_emissivity}(\varepsilon_{abs,ir}, \varepsilon_{gl,ir}) =
1/(1/0.05 + 1/0.95 - 1) \approx 0.0508$ (absorber$\to$glass, same parallel-plate formula as
$\varepsilon_{gc}$ above); $\varepsilon_{ga} = \varepsilon_{gl,ir} = 0.95$ directly (glass radiates
to ambient/sky, not a cavity). Case~1 ($\varepsilon_{abs,ir}=\varepsilon_{gl,ir}=1.0$) collapses
both to $1.0$, reproducing the blackbody base model exactly; Case~3 (both $0.0$) collapses both
to $0.0$ (no net IR exchange).
\item Absorption: $Q_{solar}=0$; desorption: $Q_{solar}\ge 0$ from the weather profile (or a
fixed baseline value for controlled sweeps).
\end{itemize}

\section{Condenser energy balance (Wilson Eq.~2, transient)}

During desorption only:
\begin{equation}
(\rho c_p L)_{cond}\,\frac{dT_{cond}}{dt}
= h_{conv,g}(T_g - T_{cond})
+ \dot{m}_{des}\,h_{fg}
+ \varepsilon_{gc}\sigma(T_g^4 - T_{cond}^4)
- A_r\,h_{amb}(T_{cond} - T_{amb})
\label{eq:condenser}
\end{equation}

where $A_r$ is the fin area ratio and $h_{conv,cond} = A_r\,h_{amb}$ (or $A_r\,h_{amb,cond}$ when a
separate fan-forced condenser convection coefficient is supplied).

\section{Mass transport (Wilson Eqs.~5--6)}

\begin{equation}
\frac{dc_w}{dt}
= \frac{g}{H_0}\,\frac{P_{sat}(T_g)}{RT_g}\,(C_R - a_w)
\label{eq:dcw}
\end{equation}

\begin{equation}
\frac{dH}{dt}
= g\,\frac{MW_{w}}{\rho_{sol}}\,\frac{P_{sat}(T_g)}{RT_g}\,(C_R - a_w)
\label{eq:dH}
\end{equation}
Wilson's Eq.~6 as printed omits the $g$ prefactor, but the bare product
$\frac{MW_w}{\rho_{sol}}\frac{P_{sat}}{RT_g}$ is dimensionless (m$^3$/mol $\times$ mol/m$^3$), so
$dH/dt$ as printed cannot carry units of m/s. Including $g$ (same mass-transfer coefficient as
Eq.~\eqref{eq:dcw}) is required both for dimensional consistency and to match the desorption
mass-rate identity $\dot{m}_{des} = -\rho_{sol}A_c\,dH/dt$ used elsewhere in the paper against the
flux-based $\dot{m}_{des} = MW_w\,g\,\frac{P_{sat}}{RT_g}(C_R-a_w)$. This implementation therefore
keeps the $g$ factor.

\textbf{Concentration ratio $C_R$:}
\begin{align}
C_R &= RH_{amb} && \text{(absorption, open chamber)} \\
C_R &= \frac{P_{sat}(T_{cond})}{P_{sat}(T_g)}\,\frac{T_g + 273.15}{T_{cond} + 273.15}
&& \text{(desorption, vapor gap)}
\end{align}

\textbf{Mass-transfer coefficient $g$:}
\begin{align}
g &= g_{chamber} && \text{(absorption)} \\
g &= h_{conv,g}\,\frac{D_{air}}{k_{air}} && \text{(desorption; Lewis-number analogy, Note~S1 Eq.~S5)}
\end{align}

\textbf{Water activity $a_w$ (PAM--LiCl):} brine activity from salt thermodynamics, capped by the
measured DVS isotherm during absorption (Note~S2). During desorption, brine activity alone.

\textbf{Saturation vapor pressure $P_{sat}(T)$:} Conde (2004) Saul--Wagner correlation, with a
Tetens-equation fallback if the primary correlation is non-finite
(\code{salt\_properties.py::saturation\_vapor\_pressure\_pa}).

\textbf{Desorption mass flux from gel inventory:}
\begin{equation}
\dot{m}_{des} = MW_w\,\bigl(-\dot{c}_w\,H - c_w\,\dot{H}\bigr), \qquad \dot{m}_{des} \ge 0
\label{eq:mdes}
\end{equation}

\section{Thermal--mass coupling constraint}

During desorption, $\dot{m}_{des}$ in Eq.~\eqref{eq:gel} must equal the flux implied by
Eqs.~\eqref{eq:dcw}--\eqref{eq:mdes}. The implementation finds $\dot{m}_{des}$ as the root of
$f(\dot{m}_{des}) = \dot{m}_{calc}(\dot{m}_{des}) - \dot{m}_{des} = 0$ (Brent's method), ensuring
Eqs.~1 and 5--6 are self-consistent at each ODE evaluation.

\section{Constraints and cycle-level outputs}

\begin{itemize}
\item $c_w \in [c_{w,\min}, c_{w,\max}]$; rates clipped at bounds.
\item $H \ge H_0$; swelling allowed in absorption ($\dot{H}\ge 0$ at $H=H_0$); no shrinkage below
$H_0$ in desorption.
\item Desorption: enforce $\dot{c}_w \le 0$, $\dot{H} \le 0$.
\item Vapor gap below \SI{7}{\milli\metre} (Wilson's thermobuoyancy/transport floor,
$L_g - H$) is treated as mass-transfer-inhibited: $g \to 0$ rather than extrapolating the
Hollands correlation outside its validated range.
\item Initial $c_w$ from fabrication equilibrium RH (DVS isotherm at $H_0$) unless overridden.
\item Daily water yield is the condenser-collected mass over one desorption phase; thermal
efficiency $\eta_{th} = \dot{m}_{water}\,h_{fg} / \int Q_{solar}\,dt$, integrated over the
desorption phase only (blending in the near-zero-solar absorption phase would dilute the ratio
without adding real signal).
\end{itemize}

\section{Physics model parameters}

Table~\ref{tab:phys-swept} lists parameters actively exposed to parameter sweeps / Bayesian
optimization, with the ranges the codebase already uses. Table~\ref{tab:phys-fixed} lists fixed
material and correlation constants with no documented sweep range.

{\small\begin{longtable}{P{3.3cm}P{1.6cm}P{2.6cm}P{3.0cm}P{3.6cm}}
\caption{Swept / design physics parameters and their documented ranges.
\label{tab:phys-swept}} \\
\toprule
\textbf{Parameter} & \textbf{Symbol} & \textbf{Baseline} & \textbf{Sweep range} & \textbf{Source} \\
\midrule
\endfirsthead
\toprule
\textbf{Parameter} & \textbf{Symbol} & \textbf{Baseline} & \textbf{Sweep range} & \textbf{Source} \\
\midrule
\endhead
\bottomrule
\endfoot
Hydrogel reference thickness & $H_0$ & \SI{4}{\milli\metre} & \SIrange{1.0}{10.0}{\milli\metre} &
\code{table\_s3.H0\_M}; range: \code{parameter\_sweep.py} / \code{lcow\_economic\_params.csv}
(\code{hydrogel\_thickness\_min/max\_m}) \\
Vapor gap & $L_g$ & \SI{40}{\milli\metre} & \SIrange{7.0}{60.0}{\milli\metre} &
\code{table\_s3.L\_G\_M}; range: \code{parameter\_sweep.py} / \code{design\_space.py}
\code{DesignBounds}; lower bound = \code{table\_s3.VAPOR\_GAP\_TRANSPORT\_MIN\_M} \\
Insulation gap & $L_{ins}$ & \SI{5}{\milli\metre} & \SIrange{1.0}{20.0}{\milli\metre} &
\code{table\_s3.L\_INS\_M}; range: \code{parameter\_sweep.py} / \code{design\_space.py} \\
Condenser fin area ratio & $A_r$ & 7.1 & \numrange{3.0}{12.0} &
\code{table\_s3.FIN\_AREA\_RATIO}; range: \code{parameter\_sweep.py} /
\code{grid\_param\_sweep.py} / \code{design\_space.py} \\
Tilt angle & $\theta$ & \SI{30}{\degree} (Fig.~2 baseline; \SI{35}{\degree} used as the Sherlock
grid-sweep baseline) & \SIrange{0}{60}{\degree} & \code{table\_s3.TILT\_DEG}; range:
\code{parameter\_sweep.py} / \code{design\_space.py} \\
Salt:polymer ratio & $S/L$ & 4.0 & \numrange{1.0}{8.0} & \code{table\_s3} default /
\code{DeviceConfig}; range: \code{parameter\_sweep.py} / \code{design\_space.py} \\
Absorber emissivity & $\varepsilon_{abs}$ & 0.95 & \numrange{0.85}{0.95} (3 discrete levels) &
\code{table\_s3.EPS\_ABS}; range: \code{sherlock\_param\_sweep.tex} /
\code{grid\_param\_sweep.py --eps-abs} \\
Glass transmittance & $\tau_{glass}$ & 0.90 & \numrange{0.80}{0.90} (3 discrete levels) &
\code{table\_s3.TAU\_GLASS}; range: \code{sherlock\_param\_sweep.tex} /
\code{grid\_param\_sweep.py --tau-glass} \\
Desorption enthalpy (LiCl) & $h_{des}$ & \SI{2.32e6}{\joule\per\kilogram} (Table~S3 COMSOL value) &
\SIrange{1.8e6}{3.2e6}{\joule\per\kilogram} & \code{table\_s3.H\_DES\_J\_PER\_KG}; range:
\code{parameter\_sweep.py}; cf. literature range in \code{salt\_heat\_of\_desorption.csv}
(LiCl $\approx$\SI{2.85e6}{\joule\per\kilogram}, D\'{i}az-Mar\'{i}n) \\
Salt weight factor (DVS-cap MW scaling, dimensionless) & --- & 1.0 & \numrange{0.826}{1.180}
(i.e.\ effective LiCl formula weight \SIrange{35}{50}{\gram\per\mole}) &
\code{parameter\_sweep.py::make\_sweep\_params} \\
Ambient convection coefficient (baseline profile) & $h_{amb}$ & \SI{10.0}{\watt\per\square\metre\per\kelvin} &
\SIrange{1.0}{12.5}{\watt\per\square\metre\per\kelvin} &
\code{weather/profiles.py::FIXED\_H\_AMB\_W\_M2\_K}; range: \code{parameter\_sweep.py} \\
Uptake RH / ambient relative humidity (baseline profile) & $RH_{amb}$ & 0.5 &
\numrange{0.15}{0.80} & \code{parameter\_sweep.py} (\code{humidity\_high}; baseline-mode
profile only) \\
Solar irradiance (baseline profile) & $Q_{solar}$ & \SI{600}{\watt\per\square\metre} &
\SIrange{400}{800}{\watt\per\square\metre} & \code{parameter\_sweep.py}
(\code{solar\_irradiance\_w\_per\_m2}; baseline-mode profile only -- real-weather mode uses
fetched irradiance instead) \\
Ambient temperature (baseline profile) & $T_{amb}$ & \SI{25}{\celsius} &
\SIrange{10}{40}{\celsius} & \code{parameter\_sweep.py} (\code{temperature\_c};
baseline-mode profile only) \\
Absorber/glass real IR emissivity pair (base case; Case~2, "selective surface") &
$\varepsilon_{abs,ir}, \varepsilon_{gl,ir}$ & (0.05, 0.95) -- \code{DeviceConfig.thermal\_params()}
default as of 2026-07 & categorical, not a continuous sweep: \code{case1} = (1.0, 1.0) reproduces
Wilson's blackbody physics exactly; \code{case3} = (0.0, 0.0) idealized optical limits &
\code{device\_config.py::DeviceConfig.thermal\_params()}; case registry:
\code{sawh\_bayesopt.design\_space.py::CASE\_EPS\_IR} -- \textbf{TBD} literature basis for the
case2/case3 numeric pairs \\
\end{longtable}}

{\small\begin{longtable}{P{3.1cm}P{1.6cm}P{3.6cm}P{4.6cm}}
\caption{Fixed physics material / correlation constants (no documented sweep range).
\label{tab:phys-fixed}} \\
\toprule
\textbf{Parameter} & \textbf{Symbol} & \textbf{Value} & \textbf{Source} \\
\midrule
\endfirsthead
\toprule
\textbf{Parameter} & \textbf{Symbol} & \textbf{Value} & \textbf{Source} \\
\midrule
\endhead
\bottomrule
\endfoot
Chamber convection (absorption) & $g_{chamber}$ & \SI{0.0085}{\metre\per\second} &
\code{table\_s3.G\_CHAMBER\_M\_S}; duplicated in economics CSV as
\code{mass\_transfer\_convection\_coefficient\_m\_s} -- \textbf{TBD} calibration source \\
Solution/brine density & $\rho_{sol}$ & \SI{1250.2}{\kilogram\per\cubic\metre} (LiCl); salt-dependent
& \code{salt\_catalog.csv} (NaCl/CaCl$_2$/MgCl$_2$ $\approx$\SI{1200}{\kilogram\per\cubic\metre}) \\
Condensation enthalpy & $h_{fg}$ & \SI{2.256e6}{\joule\per\kilogram} & \code{table\_s3.H\_FG\_J\_PER\_KG} --
\textbf{TBD} source (steam-table value assumed) \\
Gel emissivity & $\varepsilon_{gel}$ & 1.0 & \code{table\_s3.EPS\_GEL} -- \textbf{TBD} \\
Condenser (Al) emissivity & $\varepsilon_{Al}$ & 0.05 & \code{table\_s3.EPS\_AL} -- \textbf{TBD} \\
Sky temperature depression & $\Delta T_{sky}$ & \SI{0}{\kelvin} (no depression; non-zero used
only to calibrate the COMSOL Atacama field-test curves) &
\code{DeviceThermalParams.sky\_temp\_depression\_c}
-- \textbf{TBD} calibrated value / general range \\
Air thermal conductivity & $k_{air}$ & \SI{0.0286}{\watt\per\metre\per\kelvin} & \code{table\_s3.K\_AIR\_W\_M\_K} \\
Aluminum thermal conductivity & $k_{Al}$ & \SI{167.0}{\watt\per\metre\per\kelvin} & \code{table\_s3.K\_AL\_W\_M\_K} (Table~S3) \\
Hydrogel thermal conductivity & $k_{gel}$ & \SI{0.6}{\watt\per\metre\per\kelvin} & \code{table\_s3.K\_GEL\_W\_M\_K} (Table~S3) \\
Cover glass thermal conductivity & $k_{glass}$ & \SI{1.2}{\watt\per\metre\per\kelvin} & \code{table\_s3.K\_GLASS\_W\_M\_K} -- \textbf{TBD} \\
Silicone bonding-layer conductivity & $k_{silicone}$ & \SI{0.2}{\watt\per\metre\per\kelvin} & \code{table\_s3.K\_SILICONE\_W\_M\_K} -- \textbf{TBD} \\
Aluminum density / specific heat & $\rho_{Al}, c_{p,Al}$ & \SI{2700}{\kilogram\per\cubic\metre},
\SI{900}{\joule\per\kilogram\per\kelvin} & \code{table\_s3.RHO\_AL\_KG\_M3} /
\code{CP\_AL\_J\_KG\_K} -- standard aluminum handbook values, \textbf{TBD} exact citation \\
Cover-glass density / specific heat & $\rho_{glass}, c_{p,glass}$ &
\SI{2230}{\kilogram\per\cubic\metre}, \SI{830}{\joule\per\kilogram\per\kelvin} &
\code{table\_s3.RHO\_GLASS\_KG\_M3} / \code{CP\_GLASS\_J\_KG\_K} (borosilicate) -- \textbf{TBD} \\
Gel specific heat & $c_{p,gel}$ & \SI{3500}{\joule\per\kilogram\per\kelvin} &
\code{table\_s3.CP\_GEL\_J\_KG\_K} (hydrated PAM--LiCl, water-dominated) -- \textbf{TBD} \\
Water-vapor-in-air diffusivity & $D_{air}$ & \SI{2.62e-5}{\square\metre\per\second} &
\code{correlations.py::D\_AIR\_M2\_S} (Note~S1 Sh~=~Nu analogy) -- \textbf{TBD} exact source \\
Air kinematic viscosity / Prandtl number & $\nu_{air}, Pr_{air}$ & \SI{1.5e-5}{\square\metre\per\second}, 0.71
& \code{correlations.py} -- standard air properties near \SI{25}{\celsius}, \textbf{TBD} citation \\
Gravitational acceleration & $g$ & \SI{9.81}{\metre\per\square\second} & \code{correlations.py::GRAVITY\_M\_S2} \\
Stefan--Boltzmann constant & $\sigma$ & \SI{5.670374419e-8}{\watt\per\square\metre\per\kelvin\tothe{4}} & CODATA \\
Vapor-gap transport floor & --- & \SI{7}{\milli\metre} & \code{table\_s3.VAPOR\_GAP\_TRANSPORT\_MIN\_M}
(Wilson \S 2.2 thermobuoyancy/transport floor) \\
\end{longtable}}

\section{Economics model assumptions}

\begin{itemize}
\item \textbf{LCOW} is a standard fixed-charge-rate techno-economic calculation: annualized cost
divided by annual usable water volume. CAPEX is amortized via a capital recovery factor (CRF) at
a fixed discount rate over a fixed device lifetime -- not a full multi-year cash-flow model with
price escalation.
\item \textbf{NPV / payback} (\code{npv.py}) reuses the same cost components as LCOW, but
restructures CAPEX as an upfront year-0 payment against a level (non-escalating) annual net cash
flow (revenue at an assumed water price, minus OPEX), discounted at the same rate -- a
level-annuity NPV, not a full pro forma.
\item \textbf{CAPEX} is a fixed per-footprint bill of materials (BOM), summed from six line items,
scaled uniformly by a \code{total\_investment\_factor} (contingency / installation multiplier
that is easy to sweep even though the underlying line items themselves are not).
\item \textbf{Maintenance} is a fixed fraction of CAPEX per year -- not tied to any modeled
failure or service-call process.
\item \textbf{Sorbent replacement:} the hydrogel is replaced on a fixed interval
(\code{hydrogel\_lifetime\_years}); its replacement cost is the dry composite mass per footprint
area times a blended cost (catalog salt price + fixed acrylamide/additive costs, Table~S1). MOF
sorbent instead uses mass $\times$ price / lifetime directly.
\item \textbf{Utilization factor} is a single multiplicative derating of the physics-simulated
gross annual yield (daily yield $\times$ 365 $\times$ cycles/day), meant to stand in for real-world
downtime / non-ideal weather years -- not a modeled reliability/availability process.
\item \textbf{Electricity cost} only applies to the optional active/electric-heat-assist mode; the
reproduced passive-solar baseline has zero electric heat and therefore zero electricity cost.
\item \textbf{Out of scope:} water transport/distribution, land, installation labor, permitting,
and end-of-life salvage value are not modeled.
\item \textbf{No escalation:} all costs are unescalated real USD (no inflation or fuel/material
price growth curves).
\item \textbf{Water price} (NPV/payback only) is an external scenario input, not derived from any
market model.
\end{itemize}

\section{Economics governing equations}

\textbf{Capital recovery factor:}
\begin{equation}
CRF = \frac{i(1+i)^L}{(1+i)^L - 1}
\label{eq:crf}
\end{equation}
where $i$ is the discount rate and $L$ the device lifetime (years).

\textbf{Sorbent (hydrogel) replacement, annualized:}
\begin{equation}
C_{sorbent} = \frac{c_{gel}\,m_{dry}}{L_{gel}}, \qquad
c_{gel} = \frac{c_{salt}\,(S/L) + c_{acrylamide}}{1 + S/L} + c_{additives}
\label{eq:sorbent}
\end{equation}
where $m_{dry} = H_0\,\rho_{gel}$ is the dry composite mass per footprint area, $S/L$ is the
salt:polymer ratio, and $L_{gel}$ is the hydrogel replacement interval (years). (MOF sorbent:
$C_{sorbent} = m_{MOF}\,c_{MOF} / L_{gel}$ directly.)

\textbf{Annualized device cost:}
\begin{equation}
C_{annual} = CRF \cdot f_{inv}\,C_{device} + C_{sorbent} + f_{maint}\,f_{inv}\,C_{device}
+ C_{energy,fixed} + C_{elec} + C_{cycle,extra}
\label{eq:annual-cost}
\end{equation}
where $C_{device} = \sum_j C_{BOM,j}$ (six BOM line items, Table~\ref{tab:econ-fixed}),
$f_{inv}$ is the total investment factor, $f_{maint}$ the maintenance cost fraction,
$C_{elec} = p_{elec}\,\dot{Q}_{elec}\,t_{des}\cdot 365 / 1000$ is the optional active-heating
electricity cost (USD/kWh $\times$ W/m$^2$ $\times$ h/day $\times$ days/yr, converted to kWh), and
$C_{cycle,extra}$ charges a fixed energy cost per half-cycle beyond one per day.

\textbf{LCOW:}
\begin{equation}
LCOW = \frac{C_{annual}}{f_{util}\left(\dfrac{365\cdot n_{cyc}\cdot Y_{daily}}{\rho_w}\right)}
\label{eq:lcow}
\end{equation}
where $Y_{daily}$ is the simulated daily water yield (kg/m$^2$), $n_{cyc}$ the number of
half-cycles per day, $\rho_w$ the density of water (kg/m$^3$), and $f_{util}$ the utilization
factor. The cost breakdown (\code{lcow\_cost\_breakdown\_from\_daily\_yield}) reports each term
of Eq.~\eqref{eq:annual-cost} individually, divided by the same denominator, so CAPEX / sorbent /
maintenance / energy segments sum exactly to $LCOW$.

\textbf{NPV and payback:}
\begin{align}
NPV &= -f_{inv}\,C_{device} + \left(R_{annual} - OPEX_{annual}\right)\cdot
\frac{1-(1+i)^{-L}}{i}
\label{eq:npv} \\
R_{annual} &= f_{util}\left(\frac{365\cdot n_{cyc}\cdot Y_{daily}}{\rho_w}\right) p_{water}
\label{eq:revenue}
\end{align}
where $p_{water}$ is the assumed water sale price (USD/m$^3$) and $OPEX_{annual}$ is
Eq.~\eqref{eq:annual-cost} without the CAPEX/CRF term. Simple payback is
$CAPEX / (R_{annual}-OPEX_{annual})$; discounted payback solves
$1-\left(1 - i\,CAPEX/(R_{annual}-OPEX_{annual})\right) = (1+i)^{-t}$ for $t$
(\code{npv.py::npv\_from\_daily\_yield}).

\section{Economics model parameters}

Table~\ref{tab:econ-swept} lists economics parameters exposed to sweeps
(\code{parameter\_sweep.py::make\_sweep\_params}); Table~\ref{tab:econ-fixed} lists fixed
inputs (BOM line items, per-salt prices, and scenario constants) with no documented sweep range.

{\small\begin{longtable}{P{3.1cm}P{1.4cm}P{2.8cm}P{3.4cm}P{3.4cm}}
\caption{Swept economics parameters and their documented ranges.
\label{tab:econ-swept}} \\
\toprule
\textbf{Parameter} & \textbf{Symbol} & \textbf{Baseline} & \textbf{Sweep range} & \textbf{Source} \\
\midrule
\endfirsthead
\toprule
\textbf{Parameter} & \textbf{Symbol} & \textbf{Baseline} & \textbf{Sweep range} & \textbf{Source} \\
\midrule
\endhead
\bottomrule
\endfoot
Discount rate & $i$ & 0.08 & \numrange{0.04}{0.12} &
\code{lcow\_economic\_params.csv}; range: \code{parameter\_sweep.py} \\
Device lifetime & $L$ & \SI{20}{\year} & \SIrange{10}{30}{\year} (integer) &
\code{lcow\_economic\_params.csv}; range: \code{parameter\_sweep.py} \\
Utilization factor & $f_{util}$ & 0.9 & \numrange{0.7}{1.0} &
\code{lcow\_economic\_params.csv}; range: \code{parameter\_sweep.py} \\
Total investment factor & $f_{inv}$ & 1.0 & \numrange{0.7}{1.5} &
\code{lcow\_economic\_params.csv}; range: \code{parameter\_sweep.py} \\
Maintenance cost fraction & $f_{maint}$ & 0.05 \si{\per\year} & \numrange{0.02}{0.10}
\si{\per\year} & \code{lcow\_economic\_params.csv}; range: \code{parameter\_sweep.py} \\
Electricity price & $p_{elec}$ & \SI{0.10}{\usd\per\kilo\watt\hour} & \SIrange{0.05}{0.30}{\usd\per\kilo\watt\hour} &
\code{lcow\_economic\_params.csv}; range: \code{parameter\_sweep.py} \\
Acrylamide price & $c_{acrylamide}$ & \SI{1.10}{\usd\per\kilogram} &
\SIrange{0.55}{1.65}{\usd\per\kilogram} ($\pm$50\% of baseline) &
\code{lcow\_economic\_params.csv} (Table~S1 note); range: \code{parameter\_sweep.py} \\
Additives price (composite) & $c_{additives}$ & \SI{0.0072148}{\usd\per\kilogram} &
\SIrange{0.0036074}{0.0108222}{\usd\per\kilogram} ($\pm$50\% of baseline) &
\code{lcow\_economic\_params.csv} (APS, MBAA, TEMED per Table~S1); range: \code{parameter\_sweep.py} \\
Hydrogel lifetime & $L_{gel}$ & \SI{1.0}{\year} & \SIrange{0.5}{2.0}{\year} &
\code{lcow\_economic\_params.csv}; range: \code{parameter\_sweep.py} \\
Water price (NPV/payback scenario input) & $p_{water}$ & \SI{5.0}{\usd\per\cubic\metre} &
\SIrange{1.0}{25.0}{\usd\per\cubic\metre} & \code{parameter\_sweep.py}
(\code{\_BASELINE\_WATER\_PRICE\_USD\_PER\_M3}) -- \textbf{TBD} market basis \\
Salt:polymer ratio (dual physics/cost role -- see Table~\ref{tab:phys-swept}) & $S/L$ & 4.0 &
\numrange{1.0}{8.0} & See Table~\ref{tab:phys-swept}; also sets $C_{sorbent}$ via Eq.~\eqref{eq:sorbent} \\
\end{longtable}}

{\small\begin{longtable}{P{3.6cm}P{2.2cm}P{3.0cm}P{4.6cm}}
\caption{Fixed economics parameters and BOM line items (no documented sweep range).
\label{tab:econ-fixed}} \\
\toprule
\textbf{Parameter} & \textbf{Symbol} & \textbf{Value} & \textbf{Source} \\
\midrule
\endfirsthead
\toprule
\textbf{Parameter} & \textbf{Symbol} & \textbf{Value} & \textbf{Source} \\
\midrule
\endhead
\bottomrule
\endfoot
BOM: aluminum heater & $C_{BOM,1}$ & \SI{11.382}{\usd\per\square\metre} &
\code{lcow\_economic\_params.csv} (\code{device\_bom\_aluminum\_heater}) -- \textbf{TBD} itemized cost basis \\
BOM: condenser & $C_{BOM,2}$ & \SI{28.455}{\usd\per\square\metre} &
\code{lcow\_economic\_params.csv} (\code{device\_bom\_condenser}) -- \textbf{TBD} \\
BOM: acrylic enclosure & $C_{BOM,3}$ & \SI{5.95}{\usd\per\square\metre} &
\code{lcow\_economic\_params.csv} (\code{device\_bom\_acrylic\_enclosure}) -- \textbf{TBD} \\
BOM: insulation & $C_{BOM,4}$ & \SI{2.0}{\usd\per\square\metre} &
\code{lcow\_economic\_params.csv} (\code{device\_bom\_insulation}) -- \textbf{TBD} \\
BOM: clamps & $C_{BOM,5}$ & \SI{3.474}{\usd\per\square\metre} &
\code{lcow\_economic\_params.csv} (\code{device\_bom\_clamps}) -- \textbf{TBD} \\
BOM: fixtures & $C_{BOM,6}$ & \SI{1.0}{\usd\per\square\metre} &
\code{lcow\_economic\_params.csv} (\code{device\_bom\_fixtures}) -- \textbf{TBD} \\
\textit{Total device BOM} & $C_{device}$ & \textit{\SI{52.261}{\usd\per\square\metre}} (sum of
the above) & derived \\
Salt price (per catalog) & $c_{salt}$ & LiCl \num{0.55}; NaCl \num{0.045}; CaCl$_2$
\num{0.17}; MgCl$_2$ \num{0.01985} (\si{\usd\per\kilogram}) & \code{salt\_catalog.csv} --
\textbf{TBD} vendor/date basis \\
Fixed annual energy cost & $C_{energy,fixed}$ & \SI{0.0}{\usd\per\year} &
\code{lcow\_economic\_params.csv} (\code{energy\_cost\_usd\_per\_year}) -- no sweep documented \\
Extra half-cycle energy cost & --- & \SI{0.0}{\usd\per\day} &
\code{lcow\_economic\_params.csv} (\code{energy\_cost\_usd\_per\_extra\_half\_cycle\_per\_day})
-- no sweep documented \\
Desorption hours per day & $t_{des}$ & \SI{12.0}{\hour} & \code{lcow\_economic\_params.csv}
(\code{desorption\_hours\_per\_day}); tie to the physics day/night split
(\code{weather/profiles.py::PHASE\_HOURS}) is manual, not automatic -- keep the two consistent
when sweeping cycle length \\
Max electric heat (optimization bound, not baseline) & $\dot{Q}_{elec,max}$ &
\SI{1500.0}{\watt\per\square\metre} & \code{lcow\_economic\_params.csv}
(\code{max\_electric\_heat\_w\_per\_m2}) -- an optimizer bound, not a swept axis; baseline
passive device has $\dot{Q}_{elec}=0$ \\
Include desorption enthalpy in $T_g$ solve & --- & \code{true} (boolean, not swept) &
\code{lcow\_economic\_params.csv} (\code{include\_desorption\_enthalpy}) \\
Water density & $\rho_w$ & \SI{1000}{\kilogram\per\cubic\metre} &
\code{lcow\_economic\_params.csv} (\code{water\_density\_kg\_per\_l} $\times$
\code{l\_per\_m3}) \\
\end{longtable}}

\section{Integration}

The coupled physics system is integrated with SciPy \code{solve\_ivp} (Radau, stiff). State
during desorption: $\mathbf{y} = [c_w,\, H,\, T_{cond}]$; during absorption: $\mathbf{y} =
[c_w,\, H]$. LCOW/NPV are evaluated post-hoc from the resulting daily yield -- they do not feed
back into the ODE integration.

\section{Open items}

\begin{itemize}
\item Every \textbf{TBD} cell above is an as-implemented constant or sweep bound with no
literature/vendor citation currently attached in the repository -- filling these in (or replacing
the value where a better source disagrees) is the natural next step for anyone extending this
document.
\item The physics/economics duplication of the mass-transfer convection coefficient
($g_{chamber}$ in \code{table\_s3.py} vs.\ \code{mass\_transfer\_convection\_coefficient\_m\_s}
in \code{lcow\_economic\_params.csv}) is a single physical quantity kept in two files; a change
to one without the other would silently desynchronize the model.
\item \code{desorption\_hours\_per\_day} (economics, fixed at 12h) is not derived from the
physics day/night split, which can vary by site/season in real-weather mode -- a mismatch here
biases the electricity-cost term for the optional active-heating case.
\item \textbf{Case~2 default (2026-07):} \code{sawh\_bayesopt} (\code{design\_space.py},
\code{evaluator.py}, \code{bayesopt.py}, and the \code{run\_bayesopt.py}/\code{hp\_sweep.py} CLI
flags) also now defaults its own case-selection axis to \code{case2}, and \code{case1} there is
now built with an explicit \code{eps\_abs\_ir=eps\_glass\_ir=1.0} override rather than relying on
omission -- numerically identical to the old \code{None, None} encoding (verified;
$\varepsilon_{ag}=\varepsilon_{ga}=1.0$ either way), so historical \code{case1} \code{cache.jsonl}
entries remain valid, but every \code{to\_device\_config\_kwargs} call now always includes a
\code{"thermal"} key (previously omitted for \code{case1}), which changes that request's cache
key. \code{scripts/grid\_param\_sweep.py}'s \code{--eps-abs-ir}/\code{--eps-glass-ir} CLI flags
(used for the already-executed Sherlock global-grid sweep, \code{sherlock\_param\_sweep.tex})
were deliberately \emph{not} changed -- they still default to \code{None} (Case~1) to keep that
documented, already-run production sweep's baseline meaning intact; a rerun under Case~2 would
need those flags set explicitly.
\end{itemize}

\end{document}
